%% Die Sprachdatei von prag für en-GB (#377). Was eine Sprachdatei ist und
%% welche Regeln für sie gelten, steht im Kopf von helsinki-de-de.tex.

%% ---------------------------------------------------------------- Zahlen
%% Dezimalpunkt und Komma als Tausendertrenner, das € vor der Zahl und das
%% Minus davor, wie in helsinki-en-gb.tex, wo auch steht, warum \betrag das
%% Vorzeichen selbst setzt.
\sisetup{
  output-decimal-marker = {.},
  group-separator       = {,},
}
\makeatletter
\newcommand{\betrag}[1]{\betrag@zeichen#1\relax}
\def\betrag@zeichen#1#2\relax{\ifx-#1\textminus€\num{#2}\else€\num{#1#2}\fi}
\makeatother
\newcommand{\prozent}[1]{\num{#1}\%}

%% ------------------------------------------------------------- Seiten
\newcommand{\wortseitevon}[2]{Page #1 of #2}
\newcommand{\wortfortsetzung}{Continued on the next page}

%% ------------------------------------------------------------ Kopfangaben
\newcommand{\wortrechnungsnr}{Invoice no.:}
\newcommand{\wortkundennr}{Customer no.:}
\newcommand{\wortauftragsnr}{Our order no.:}
\newcommand{\wortdatum}{Date:}
\newcommand{\wortvereinnahmtam}{Received on:}
\newcommand{\wortreferenz}{Your reference:}

%% ----------------------------------------------------- Positionen, Summen
%% Kurze Wörter im Spaltenkopf: Die Spalten des Gitters sind so schmal wie
%% ihr deutscher Kopf, "Quantity" passte nicht über die Menge.
\newcommand{\wortpos}{No.}
\newcommand{\wortbezeichnung}{Item}
\newcommand{\wortmenge}{Qty}
\newcommand{\worteinheit}{Unit}
\newcommand{\worteinzelpreis}{Unit price}
\newcommand{\wortrabatt}{Discount}
\newcommand{\wortgesamt}{Total}
\newcommand{\wortzwischensumme}{Subtotal (net)}
%% Satz und, bei mehreren Steuergruppen, ihre Grundlage – sonst ist sie leer.
\newcommand{\wortsteuer}[2]{VAT #1\ifblank{#2}{}{ on #2}}
\newcommand{\wortgesamtbetrag}{Total amount}

%% --------------------------------------------------------------------- Fuß
\newcommand{\wortustid}{VAT ID}
\newcommand{\wortsteuernummer}{Tax number}
\newcommand{\wortregistergericht}{Registry court}
\newcommand{\wortgeschaeftsfuehrung}{Managing director}
